\documentclass[a4paper,11pt]{article}

\usepackage[utf8]{inputenc}
\usepackage[T1]{fontenc}
\usepackage{geometry}
\usepackage{booktabs}
\usepackage{array}
\usepackage{longtable}
\usepackage{xcolor}
\usepackage{listings}

\geometry{
    top=2cm,
    bottom=2cm,
    left=2cm,
    right=2cm
}

\setlength{\parindent}{0pt}
\setlength{\parskip}{5pt}

\lstset{
    basicstyle=\ttfamily\small,
    columns=fullflexible,
    frame=single,
    breaklines=true
}

\title{\textbf{DEOS Fault Management Tanımı}}
\author{DEOS -- Dokuz Eylül Otonom Sistemler}
\date{}

\begin{document}

\maketitle

\section{Amaç}

Bu doküman, DEOS CAN-FD ICD içerisinde tanımlanan
\texttt{DIAGNOSTIC} servisine ait fault yönetim mekanizmasını
ayrıntılandırır.

Tanımlanan mekanizma yalnızca BMS sistemi için değil,
DEOS protokolüne dahil olan tüm fiziksel ve sanal node'lar için
geçerlidir.

Her node kendi içerisinde oluşan fault kayıtlarını yerel olarak tutar.
Fault kayıtlarının diğer node'lara aktarılması
\texttt{DEOS\_SERVICE\_DIAGNOSTIC} servisi üzerinden gerçekleştirilir.

Bu doküman;

\begin{itemize}[nosep]
    \item fault kayıt yapısını,
    \item fault kimliklendirme yöntemini,
    \item \texttt{GET\_FAULTS} davranışını,
    \item fault response formatını,
    \item \texttt{CLEAR\_FAULTS} davranışını,
    \item node bazlı başlangıç fault registry'lerini
\end{itemize}

tanımlar.

\section{Diagnostic Service}

Diagnostic Service ID aşağıdaki gibidir:

\begin{center}
\begin{tabular}{c c l}
\toprule
\textbf{Service ID} & \textbf{Symbol} & \textbf{Açıklama} \\
\midrule
0x06 & DEOS\_SERVICE\_DIAGNOSTIC & Node diagnostik ve fault işlemleri \\
\bottomrule
\end{tabular}
\end{center}

Bu servis altında mevcut fault mekanizması için kullanılan komutlar:

\begin{center}
\begin{tabular}{c l l}
\toprule
\textbf{Command ID} & \textbf{Symbol} & \textbf{Amaç} \\
\midrule
0x02 & DEOS\_CMD\_DIAG\_GET\_FAULTS   & Node üzerindeki fault kayıtlarını isteme \\
0x03 & DEOS\_CMD\_DIAG\_CLEAR\_FAULTS & Node üzerindeki fault kayıtlarını temizleme \\
\bottomrule
\end{tabular}
\end{center}

\section{Fault Kayıt Modeli}

Her DEOS node'u kendi fault kayıt alanından sorumludur.

Fault kayıt alanının RAM, Flash, EEPROM veya başka bir bellek yapısında
tutulması protokol tarafından zorunlu tutulmaz.
Bu seçim ilgili node yazılımının implementasyonuna bırakılmıştır.

Her fault kaydı aşağıdaki mantıksal alanlardan oluşur:

\begin{center}
\begin{tabular}{l c l}
\toprule
\textbf{Alan} & \textbf{Tip} & \textbf{Açıklama} \\
\midrule
Fault ID &
uint16\_t &
Node içerisindeki fault kimliği \\

Severity &
uint8\_t &
Fault önem seviyesi \\

State &
uint8\_t &
Fault'un mevcut durumu \\

Occurrence Count &
uint16\_t &
Fault'un kaç kez oluştuğu \\

Last Occurrence Time &
uint32\_t &
Fault'un son oluştuğu node uptime zamanı, ms \\
\bottomrule
\end{tabular}
\end{center}

Önerilen yazılım gösterimi:

\begin{lstlisting}[language=C]
typedef struct
{
    uint16_t fault_id;
    uint8_t  severity;
    uint8_t  state;
    uint16_t occurrence_count;
    uint32_t last_occurrence_ms;
} deos_fault_record_t;
\end{lstlisting}

\section{Fault ID Namespace}

DEOS sisteminde global bir fault registry kullanılmaz.

Bir fault'un kimliği aşağıdaki ikili ile belirlenir:

\begin{center}
\texttt{Source Node ID + Fault ID}
\end{center}

Dolayısıyla;

\begin{lstlisting}
STEERING   / Fault 0x0101
\end{lstlisting}

ile

\begin{lstlisting}
TRACTION   / Fault 0x0101
\end{lstlisting}

aynı fault olmak zorunda değildir.

Her node kendi fault namespace'ine sahiptir.

Bununla birlikte yazılım bütünlüğünü artırmak amacıyla,
birçok node'da ortak olarak ortaya çıkabilecek temel yazılım ve
haberleşme fault'ları için düşük Fault ID değerlerinin aynı anlamda
kullanılması önerilir.

\subsection{Ortak Yerel Fault ID'leri}

Aşağıdaki fault ID'leri node içerisinde ihtiyaç duyulması halinde
aynı anlamla kullanılmalıdır.

\begin{center}
\begin{tabular}{c l}
\toprule
\textbf{Fault ID} & \textbf{Önerilen Anlam} \\
\midrule
0x0000 & Reserved -- End of Fault List \\
0x0001 & Internal Software Error \\
0x0002 & Watchdog Reset Detected \\
0x0003 & Invalid Configuration \\
0x0004 & Communication Timeout \\
0x0005 & RX Buffer / Queue Overflow \\
0x0006 & TX Failure \\
0x0007 & Protocol / Invalid Message Error \\
0x0008 & Local Storage Error \\
\bottomrule
\end{tabular}
\end{center}

Bu tablo global bir fault sistemi oluşturmaz.

Örneğin \texttt{0x0004} kodu Steering node'unda oluştuğunda
Steering node'un kendi communication timeout kaydını,
BMS Gateway üzerinde oluştuğunda ise BMS Gateway'in kendi
communication timeout kaydını ifade eder.

Node'a özgü fault'ların \texttt{0x0100} ve üzerindeki değerlerden
başlatılması önerilir.

\section{Severity}

Fault önem derecesi aşağıdaki şekilde tanımlanır:

\begin{center}
\begin{tabular}{c l}
\toprule
\textbf{Değer} & \textbf{Anlam} \\
\midrule
0x01 & INFO \\
0x02 & WARNING \\
0x03 & ERROR \\
0x04 & CRITICAL \\
\bottomrule
\end{tabular}
\end{center}

Severity değeri, fault'un sistem üzerindeki etkisini ifade eder.
Fault ID'nin kendisinin bir parçası değildir.

\section{Fault State}

Fault state alanı aşağıdaki şekilde tanımlanır:

\begin{center}
\begin{tabular}{c l}
\toprule
\textbf{Değer} & \textbf{Anlam} \\
\midrule
0x00 & INACTIVE \\
0x01 & ACTIVE \\
0x02 & LATCHED \\
\bottomrule
\end{tabular}
\end{center}

\texttt{ACTIVE}, fault condition'ın halen mevcut olduğunu gösterir.

\texttt{INACTIVE}, fault'un daha önce oluştuğunu ancak mevcut
condition'ın artık ortadan kalktığını gösterir.

\texttt{LATCHED}, fault condition ortadan kalksa dahi kaydın
\texttt{CLEAR\_FAULTS} işlemine kadar tutulabileceğini ifade eder.

\section{Fault Oluşturma}

Bir node içerisinde tanımlı bir hata condition'ı tespit edildiğinde
ilgili fault kaydı oluşturulur veya mevcut kayıt güncellenir.

Aynı Fault ID tekrar oluştuğunda yeni bir ayrı kayıt oluşturulması
zorunlu değildir.

Önerilen davranış:

\begin{lstlisting}
fault.state = ACTIVE;
fault.occurrence_count++;
fault.last_occurrence_ms = current_uptime_ms;
\end{lstlisting}

Fault condition ortadan kalktığında ilgili kayıt implementasyona göre
\texttt{INACTIVE} veya \texttt{LATCHED} durumuna geçirilebilir.

\section{GET\_FAULTS}

\texttt{GET\_FAULTS}, hedef node üzerinde kayıtlı fault'ların tamamını
istemek için kullanılır.

Mesaj:

\begin{lstlisting}
Message Class = DEOS_CLASS_REQUEST
Service       = DEOS_SERVICE_DIAGNOSTIC
Command       = DEOS_CMD_DIAG_GET_FAULTS
\end{lstlisting}

GET\_FAULTS request mesajında command-specific payload bulunmaz.

Genel DEOS data alanı:

\begin{center}
\begin{tabular}{c c l}
\toprule
\textbf{Byte} & \textbf{Alan} & \textbf{Değer} \\
\midrule
0 & Protocol Version & DEOS protocol version \\
1 & Sequence         & Gönderen node TX sequence \\
2 & Command          & 0x02 -- GET\_FAULTS \\
\bottomrule
\end{tabular}
\end{center}

\subsection{GET\_FAULTS Davranışı}

Bir node kendisine yöneltilmiş geçerli bir
\texttt{GET\_FAULTS} mesajı aldığında, o anda kendi fault kayıt
alanında bulunan fault'ların tamamını gönderir.

Her fault bağımsız bir CAN-FD frame'i içerisinde gönderilir.

Bir frame içerisinde birden fazla fault bulunmaz.

Ardışık fault response mesajları arasında nominal olarak:

\begin{center}
\textbf{100 ms}
\end{center}

süre bırakılır.

İlk fault response mesajının gönderilmesi için 100 ms beklenmesi
zorunlu değildir.

Örnek:

\begin{lstlisting}
MAIN_STM32                  STEERING
    |                           |
    |---- GET_FAULTS ---------->|
    |                           |
    |<--- Fault 0x0101 ---------|
    |                           |
    |       100 ms              |
    |                           |
    |<--- Fault 0x0103 ---------|
    |                           |
    |       100 ms              |
    |                           |
    |<--- Fault 0x0107 ---------|
    |                           |
    |       100 ms              |
    |                           |
    |<--- END ------------------|
\end{lstlisting}

GET\_FAULTS işlemi index, offset veya pagination mekanizması kullanmaz.

Tek bir GET\_FAULTS isteği, ilgili node üzerindeki mevcut fault
kayıtlarının tamamının aktarılmasını başlatır.

\section{Fault Response Mesajı}

Her fault aşağıdaki mesaj yapısı ile gönderilir:

\begin{lstlisting}
Message Class = DEOS_CLASS_RESPONSE
Service       = DEOS_SERVICE_DIAGNOSTIC
Command       = DEOS_CMD_DIAG_GET_FAULTS
\end{lstlisting}

Response Source Node, fault'un sahibi olan node'dur.

Response Destination Node ise GET\_FAULTS isteğini gönderen node'dur.

Fault response command payload formatı:

\begin{center}
\begin{tabular}{c c c l}
\toprule
\textbf{Payload Offset} &
\textbf{Boyut} &
\textbf{Tip} &
\textbf{Alan} \\
\midrule
0 & 2 byte & uint16\_t & Fault ID \\
2 & 1 byte & uint8\_t  & Severity \\
3 & 1 byte & uint8\_t  & State \\
4 & 2 byte & uint16\_t & Occurrence Count \\
6 & 4 byte & uint32\_t & Last Occurrence Time [ms] \\
\bottomrule
\end{tabular}
\end{center}

Tüm çok baytlı alanlar DEOS ICD genel kuralına uygun şekilde
Little-Endian olarak aktarılır.

CAN-FD data alanı genel olarak:

\begin{center}
\begin{tabular}{c l}
\toprule
\textbf{Byte} & \textbf{Alan} \\
\midrule
0 & Protocol Version \\
1 & Sequence \\
2 & GET\_FAULTS (0x02) \\
3--4 & Fault ID \\
5 & Severity \\
6 & State \\
7--8 & Occurrence Count \\
9--12 & Last Occurrence Time \\
\bottomrule
\end{tabular}
\end{center}

Semantik mesaj uzunluğu 13 byte'tır.
CAN-FD DLC nedeniyle kullanılan fiziksel frame uzunluğu daha büyük
olursa kalan byte'lar sıfır ile doldurulur.

\section{Fault Listesinin Sonu}

\texttt{Fault ID = 0x0000} gerçek bir fault olarak kullanılmaz.

Bu değer:

\begin{center}
\textbf{END\_OF\_FAULT\_LIST}
\end{center}

anlamına gelir.

Node bütün fault kayıtlarını gönderdikten sonra son olarak
\texttt{Fault ID = 0x0000} içeren bir response gönderir.

END mesajında diğer fault alanları sıfır olmalıdır:

\begin{lstlisting}
Fault ID        = 0x0000
Severity        = 0x00
State           = 0x00
Occurrence Count= 0x0000
Timestamp       = 0x00000000
\end{lstlisting}

Node üzerinde hiç fault yoksa GET\_FAULTS isteğine doğrudan
END\_OF\_FAULT\_LIST mesajı ile cevap verilir.

\section{CLEAR\_FAULTS}

Fault kayıtlarının temizlenmesi için:

\begin{lstlisting}
Message Class = DEOS_CLASS_COMMAND
Service       = DEOS_SERVICE_DIAGNOSTIC
Command       = DEOS_CMD_DIAG_CLEAR_FAULTS
\end{lstlisting}

kullanılır.

Command-specific payload bulunmaz.

Geçerli bir CLEAR\_FAULTS mesajı alındığında hedef node kendi fault
kayıt alanını temizler.

CLEAR\_FAULTS işlemi node'un gerçek çalışma koşullarını değiştirmez.

Örneğin CAN haberleşmesi halen \texttt{BUS\_OFF} durumundaysa,
CAN fault kaydı temizlendikten sonra hata condition'ı tekrar
algılanabilir ve aynı fault yeniden oluşturulabilir.

Bu nedenle CLEAR\_FAULTS:

\begin{center}
\textbf{fault condition'ı değil, saklanan fault kayıtlarını temizler.}
\end{center}

\section{Node Bazlı Fault Registry}

Aşağıdaki tablolar başlangıç fault registry önerilerini tanımlar.

Her node yalnızca kendi donanımı ve yazılımı için anlamlı fault'ları
uygulamak zorundadır.

\subsection{Main STM32}

\begin{center}
\begin{tabular}{c l}
\toprule
\textbf{ID} & \textbf{Fault} \\
\midrule
0x0100 & CAN-FD Bus Off \\
0x0101 & Ethernet Link / Communication Failure \\
0x0102 & LoRa Communication Failure \\
0x0103 & Router Queue Overflow \\
0x0104 & Invalid / Unknown Route \\
\bottomrule
\end{tabular}
\end{center}

\subsection{Traction Node}

\begin{center}
\begin{tabular}{c l}
\toprule
\textbf{ID} & \textbf{Fault} \\
\midrule
0x0100 & Motor / Inverter Communication Lost \\
0x0101 & Motor Overcurrent \\
0x0102 & Motor Overtemperature \\
0x0103 & Speed Feedback Invalid \\
0x0104 & Overspeed Detected \\
0x0105 & Actuator / Motor Response Timeout \\
\bottomrule
\end{tabular}
\end{center}

\subsection{Steering Node}

\begin{center}
\begin{tabular}{c l}
\toprule
\textbf{ID} & \textbf{Fault} \\
\midrule
0x0100 & Steering Encoder Communication Lost \\
0x0101 & Steering Position Invalid \\
0x0102 & Steering Angle Out of Range \\
0x0103 & Steering Motor Overcurrent \\
0x0104 & Steering Motor / Actuator Stall \\
0x0105 & Steering Response Timeout \\
\bottomrule
\end{tabular}
\end{center}

\subsection{Brake Node}

\begin{center}
\begin{tabular}{c l}
\toprule
\textbf{ID} & \textbf{Fault} \\
\midrule
0x0100 & Brake Position Feedback Invalid \\
0x0101 & Brake Actuator Overcurrent \\
0x0102 & Brake Actuator Stall \\
0x0103 & Brake Position Out of Range \\
0x0104 & Brake Response Timeout \\
\bottomrule
\end{tabular}
\end{center}

\subsection{BMS Main ve BMS Aux}

BMS_MAIN ve BMS_AUX ayrı DEOS Node ID'lerine sahip olduğundan,
aynı Fault ID registry'sini birbirlerinden bağımsız olarak
kullanabilirler.

\begin{center}
\begin{tabular}{c l}
\toprule
\textbf{ID} & \textbf{Fault} \\
\midrule
0x0100 & Pack Overvoltage \\
0x0101 & Pack Undervoltage \\
0x0102 & Charge Overcurrent \\
0x0103 & Discharge Overcurrent \\
0x0104 & Overtemperature \\
0x0105 & Undertemperature \\
0x0106 & Cell Voltage Imbalance \\
0x0107 & BMS Internal Fault \\
0x0108 & Invalid BMS Measurement \\
\bottomrule
\end{tabular}
\end{center}

BMS fault'ları gateway tarafından DALY BMS bilgilerinden üretilse
bile DEOS mesajındaki Source Node ilgili BMS node'u olmalıdır.

Örneğin BMS_MAIN fault'u:

\begin{lstlisting}
Source  = DEOS_NODE_BMS_MAIN
Service = DEOS_SERVICE_DIAGNOSTIC
\end{lstlisting}

olarak gönderilir.

\subsection{BMS Gateway}

\begin{center}
\begin{tabular}{c l}
\toprule
\textbf{ID} & \textbf{Fault} \\
\midrule
0x0100 & BMS Main Communication Lost \\
0x0101 & BMS Aux Communication Lost \\
0x0102 & Classic CAN Bus Off \\
0x0103 & CAN-FD Bus Off \\
0x0104 & Invalid BMS Frame \\
0x0105 & BMS Protocol / CRC Error \\
\bottomrule
\end{tabular}
\end{center}

Gateway'in kendi fault'larında Source Node:

\begin{lstlisting}
DEOS_NODE_BMS_GATEWAY
\end{lstlisting}

olmalıdır.

Gateway'in BMS_MAIN veya BMS_AUX adına oluşturduğu fault'larda ise
Source Node ilgili mantıksal BMS Node ID'si olarak korunmalıdır.

\subsection{Textual Node}

\begin{center}
\begin{tabular}{c l}
\toprule
\textbf{ID} & \textbf{Fault} \\
\midrule
0x0100 & DEOS Transport Connection Lost \\
0x0101 & Invalid Message Serialization / Parsing \\
0x0102 & Configuration Operation Failed \\
\bottomrule
\end{tabular}
\end{center}

\subsection{micro-ROS Node}

\begin{center}
\begin{tabular}{c l}
\toprule
\textbf{ID} & \textbf{Fault} \\
\midrule
0x0100 & ROS / micro-ROS Connection Lost \\
0x0101 & ROS Agent Unreachable \\
0x0102 & DEOS--ROS Message Conversion Failed \\
0x0103 & Required Topic / Endpoint Unavailable \\
\bottomrule
\end{tabular}
\end{center}

\subsection{Ground Control Node}

\begin{center}
\begin{tabular}{c l}
\toprule
\textbf{ID} & \textbf{Fault} \\
\midrule
0x0100 & LoRa Communication Lost \\
0x0101 & Invalid LoRa / DEOS Message \\
0x0102 & Transport Timeout \\
\bottomrule
\end{tabular}
\end{center}

\section{Örnek GET\_FAULTS İşlemi}

Main STM32, Steering node üzerinde kayıtlı fault'ları istemektedir.

Request:

\begin{lstlisting}
Source        = DEOS_NODE_MAIN_STM32
Destination   = DEOS_NODE_STEERING
Message Class = DEOS_CLASS_REQUEST
Service       = DEOS_SERVICE_DIAGNOSTIC
Command       = DEOS_CMD_DIAG_GET_FAULTS
\end{lstlisting}

Steering node üzerinde aşağıdaki fault'ların bulunduğu varsayılsın:

\begin{lstlisting}
0x0100  Steering Encoder Communication Lost
0x0103  Steering Motor Overcurrent
\end{lstlisting}

Mesaj akışı:

\begin{lstlisting}
MAIN_STM32                         STEERING
    |                                  |
    |------ GET_FAULTS --------------->|
    |                                  |
    |<----- Fault 0x0100 --------------|
    |                                  |
    |          100 ms                  |
    |                                  |
    |<----- Fault 0x0103 --------------|
    |                                  |
    |          100 ms                  |
    |                                  |
    |<----- END (0x0000) --------------|
\end{lstlisting}

Her fault response mesajının Source alanı
\texttt{DEOS\_NODE\_STEERING} olarak kalır.

\section{Örnek CLEAR\_FAULTS İşlemi}

\begin{lstlisting}
Source        = DEOS_NODE_MAIN_STM32
Destination   = DEOS_NODE_STEERING
Message Class = DEOS_CLASS_COMMAND
Service       = DEOS_SERVICE_DIAGNOSTIC
Command       = DEOS_CMD_DIAG_CLEAR_FAULTS
\end{lstlisting}

Steering node bu mesajı aldıktan sonra saklanan fault kayıtlarını
temizler.

Devam eden fiziksel veya yazılımsal bir hata mevcutsa ilgili
fault daha sonra yeniden oluşturulabilir.

\section{Özet Davranış}

Fault mekanizmasının temel kuralları aşağıdaki gibidir:

\begin{enumerate}
    \item Her node kendi fault kayıt alanına sahiptir.
    
    \item Global bir fault registry yoktur.
    
    \item Fault kimliği Source Node ID ve Fault ID'nin birlikte
    değerlendirilmesi ile belirlenir.
    
    \item Düşük Fault ID değerleri ortak yazılım ve haberleşme
    hataları için aynı anlamda kullanılabilir.
    
    \item Node'a özgü fault'ların 0x0100 ve üzerindeki ID'lerden
    başlaması önerilir.
    
    \item GET\_FAULTS tek bir request ile node üzerindeki bütün
    fault kayıtlarının gönderilmesini başlatır.
    
    \item Her CAN-FD response frame'i yalnızca bir fault taşır.
    
    \item Ardışık fault frame'leri 100 ms aralıklarla gönderilir.
    
    \item Fault listesinin bittiği Fault ID = 0x0000 ile bildirilir.
    
    \item CLEAR\_FAULTS node'un saklanan fault kayıtlarını temizler.
    
    \item Devam eden bir hata condition'ı CLEAR\_FAULTS sonrasında
    tekrar fault oluşturabilir.
\end{enumerate}

\end{document}